% SGA 3, Expose III, Loop 1 English reconstruction.
% Scope: Remark 0.3 through Corollary 0.7, stopping immediately before
% Corollary 0.7 bis.
% Authority: Polo--Gille Exp3-14oct24.pdf, component-local pages 14--17
% (combined-reader pages 124--127; printed pages 114--117). Force-OCR is a
% drafting locator; English wording was compared with Jacob Reinhold's
% CC BY 4.0 Markdown lineage. The PDF controls text, formulas, numbering,
% notes, and diagram images.

\setcounter{footnote}{20}

% Source page: Exp3 component-local p. 14; combined-reader p. 124; printed p. 114.
\medskip
\noindent\SGAInlineRefTarget{sga3:iii:remark:0-3}\textbf{注 0.3.}
假设 \(X^+\)-概形 \(Y\) 在 \(S\) 上平坦（参见 SGA~1,
IV）. 此时可写成
\[
 \operatorname{Hom}_{X^+}(Y,L_X)
 =
 \operatorname{Hom}_{\mathscr O_{Y_0}}
 \bigl(g_0^*(\Omega^1_{X_0/S_0}),
       \mathscr J\otimes_{\mathscr O_{S_0}}\mathscr O_{Y_0}\bigr).
\]

\medskip
\noindent\SGAInlineRefTarget{sga3:iii:remark:0-4}\textbf{注 0.4.}
设 \(\pi_0:X_0\to S_0\) 是结构态射, 并假设存在一个
\(\mathscr O_{S_0}\)-模 \(\omega^1_{X_0/S_0}\), 使得
\[
 \Omega^1_{X_0/S_0}\simeq
 \pi_0^*(\omega^1_{X_0/S_0}).
\]
（特别地, 当 \(X_0\) 是一个 \(S_0\)-群时就是如此;
参见 \SGARefLink{sga3:ii:subsection:4-11}{II~4.11}.）
由下式定义 \(S\) 上的函子 \(L'_X\):
\begin{equation}
 \operatorname{Hom}_S(Y,L'_X)
 =
 \operatorname{Hom}_{\mathscr O_{Y_0}}
 \bigl(\omega^1_{X_0/S_0}\otimes_{\mathscr O_{S_0}}\mathscr O_{Y_0},
       \mathscr J\mathscr O_Y\bigr).
 \tag{0.4.1}\label{sga3:iii:remark:0-4:equation:0-4-1}
\end{equation}
则对每个 \(X^+\)-概形 \(Y\), 有
\[
 \operatorname{Hom}_{X^+}(Y,L_X)=\operatorname{Hom}_S(Y,L'_X)
\]
; 换言之,
\[
 L_X=L'_X\mathbin{\times_S}X^+.
\]

由于
\(L_X\mathbin{\times_{X^+}}X=L'_X\mathbin{\times_S}X\),
\(L_X\) 在 \(X\) 上的作用诱导出 \(L'_X\) 在 \(X\) 上的作用,
而此作用保持 \(p_X:X\to X^+\).\footnote{编者注: 以下对原文作了展开.}
事实上, 若 \(Y\) 是一个 \(S\)-概形,
\(h:Y\to X\) 是一个 \(S\)-态射, 且 \(m\in L'_X(Y)\),
则 \(h\) 和 \(m\cdot h\) 在 \(Y_{\mathscr J}\) 上的限制相同;
也就是说,
\[
 p_X(m\cdot h)=p_X(h).
\]

\medskip
\noindent\SGAInlineRefTarget{sga3:iii:remark:0-5}\textbf{注 0.5.}
沿用 \SGARefLink{sga3:iii:remark:0-4}{0.4} 的假设与记号,
并再假设 \(Y\) 是 \(S\) 上平坦的 \(S\)-概形. 则
\[
 \operatorname{Hom}_{X^+}(Y,L_X)
 =\operatorname{Hom}_S(Y,L'_X)
 =\operatorname{Hom}_{S_0}(Y_0,L_{0X}),
\]
其中, 取值于阿贝尔群的 \(S_0\)-函子 \(L_{0X}\)
由下列恒等式定义（这里 \(T\) 是变动的 \(S_0\)-概形）:
\begin{equation}
 \operatorname{Hom}_{S_0}(T,L_{0X})
 =
 \operatorname{Hom}_{\mathscr O_T}
 \bigl(\omega^1_{X_0/S_0}\otimes_{\mathscr O_{S_0}}\mathscr O_T,
       \mathscr J\otimes_{\mathscr O_{S_0}}\mathscr O_T\bigr).
 \tag{0.5.1}\label{sga3:iii:remark:0-5:equation:0-5-1}
\end{equation}
所以, 使用 II~1 的记号, 我们已证明 \(L'_X\) 与
\(\prod_{S_0/S}L_{0X}\) 在 \((\mathbf{Sch})/S\) 的下述全子范畴上的限制相同:
其对象为 \(S\) 上平坦的 \(S\)-概形.

\medskip
\noindent\SGAInlineRefTarget{sga3:iii:remark:0-6}\textbf{注 0.6.}
沿用 \SGARefLink{sga3:iii:remark:0-5}{0.5} 的假设与记号,
\footnote{编者注: 新增了以下假设, 原文中仅隐含使用了这一假设.}
再假设 \(\pi_0:X_0\to S_0\) 有一个截面 \(\varepsilon_0\). 则
\[
 \omega^1_{X_0/S_0}\simeq
 \varepsilon_0^*(\Omega^1_{X_0/S_0}).
\]
首先, 不依赖于最后这一假设, 有
\[
 \operatorname{Hom}_{S_0}(T,L_{0X})
 =
 \Gamma\!\left(
   T,
   \mathscr{H}om_{\mathscr O_T}
   \bigl(\omega^1_{X_0/S_0}\otimes_{\mathscr O_{S_0}}\mathscr O_T,
         \mathscr J\otimes_{\mathscr O_{S_0}}\mathscr O_T\bigr)
 \right).
\]

% Source page: Exp3 component-local p. 15; combined-reader p. 125; printed p. 115.
现假设 \(\omega^1_{X_0/S_0}\) 有限呈示（参见 EGA~0\textsubscript{I}, 5.2.5）;
特别地, 当 \(X_0\) 在 \(S_0\) 上局部有限呈示时就是如此
（参见 EGA~IV\textsubscript{4}, 16.4.22）. 若 \(T\) 在 \(S_0\) 上平坦,
则由 EGA~0\textsubscript{I}, 6.7.6, 得
\[
 \mathscr{H}om_{\mathscr O_T}
 \bigl(\omega^1_{X_0/S_0}\otimes_{\mathscr O_{S_0}}\mathscr O_T,
       \mathscr J\otimes_{\mathscr O_{S_0}}\mathscr O_T\bigr)
 \simeq
 \mathscr{H}om_{\mathscr O_{S_0}}
 \bigl(\omega^1_{X_0/S_0},\mathscr J\bigr)
 \otimes_{\mathscr O_{S_0}}\mathscr O_T,
\]
从而
\[
 \operatorname{Hom}_{S_0}(T,L_{0X})
 =
 \Gamma\!\left(
   T,
   \mathscr{H}om_{\mathscr O_{S_0}}
   \bigl(\omega^1_{X_0/S_0},\mathscr J\bigr)
   \otimes_{\mathscr O_{S_0}}\mathscr O_T
 \right).
\]
使用 \SGARefLink{sga3:i:definition:4-6-1}{I~4.6.1} 中的记号
\(\mathbf W(\,\cdot\,)\), 我们由此证明了: 对每个在
\(S_0\) 上平坦的 \(S_0\)-概形 \(T\), 有
\[
 \operatorname{Hom}_{S_0}(T,L_{0X})
 =
 \operatorname{Hom}_{S_0}
 \left(T,
 \mathbf W\!\left(
   \mathscr{H}om_{\mathscr O_{S_0}}
   (\omega^1_{X_0/S_0},\mathscr J)
 \right)\right).
\]

总之, 若 \(\omega^1_{X_0/S_0}\) 有限呈示,
并限制到 \(S\) 上平坦的 \(S\)-概形范畴, 则
\begin{equation}
 L'_X
 =
 \prod_{S_0/S}
 \mathbf W\!\left(
   \mathscr{H}om_{\mathscr O_{S_0}}
   (\omega^1_{X_0/S_0},\mathscr J)
 \right),
 \tag{0.6.1}\label{sga3:iii:remark:0-6:equation:0-6-1}
\end{equation}
且由 EGA~I, 9.1.1,
\(\mathscr{H}om_{\mathscr O_{S_0}}
(\omega^1_{X_0/S_0},\mathscr J)\) 是拟凝聚
\(\mathscr O_{S_0}\)-模.

最后, 再假设 \(\omega^1_{X_0/S_0}\) 是有限秩局部自由的,
例如 \(X_0\) 在 \(S_0\) 上平滑（此时它自动在 \(S_0\) 上局部有限呈示）.
则
\begin{equation}
 \mathscr{H}om_{\mathscr O_{S_0}}
   (\omega^1_{X_0/S_0},\mathscr J)
 \simeq
 \underline{\operatorname{Lie}}(X_0/S_0)
 \otimes_{\mathscr O_{S_0}}\mathscr J,
 \tag{0.6.2}\label{sga3:iii:remark:0-6:equation:0-6-2}
\end{equation}
其中, 采用略不严格的说法（因为 \(X_0\) 未必是一个
\(S_0\)-群）, 以 \(\underline{\operatorname{Lie}}(X_0/S_0)\)
表示 \(\mathscr O_{S_0}\)-模 \(\omega^1_{X_0/S_0}\) 的对偶.
\footnote{编者注: 这一说法由 \SGARefLink{sga3:ii:proposition:3-3}{II~3.3}
和 \SGARefLink{sga3:iii:paragraph:4-11}{4.11} 说明.
事实上, 函子 \(L^{\varepsilon_0}_{X_0/S_0}\), 即“\(X_0\)
在 \(S_0\) 上于 \(\varepsilon_0\) 处的切空间”, 可由 \(S_0\) 上的向量丛
\(\mathbf V(\varepsilon_0^*(\Omega^1_{X_0/S_0}))
=\mathbf V(\omega^1_{X_0/S_0})\) 表示.
其截面层是对偶模
\(\mathscr{H}om_{\mathscr O_{S_0}}
(\omega^1_{X_0/S_0},\mathscr O_{S_0})\);
当 \(X_0\) 是一个 \(S_0\)-群, 且 \(\varepsilon_0\) 为单位截面时,
它就是 \(\underline{\operatorname{Lie}}(X_0/S_0)\).}

\SGARefLink{sga3:iii:proposition:0-2}{命题~0.2}（及其证明）有两个重要的推论.
\footnote{编者注: 在 \SGARefLink{sga3:iii:corollary:0-7}{0.7} 之外,
新增了在 \SGARefLink{sga3:iii:proposition:0-8}{0.8} 的证明中隐含使用的
\SGARefLink{sga3:iii:corollary:0-7-bis}{推论 0.7 bis}.
我们指出, \SGARefLink{sga3:iii:proposition:0-8}{0.8}--\SGARefLink{sga3:iii:corollary:0-12}{0.12}
以及 \SGARefLink{sga3:iii:corollary:0-17}{0.17} 都是
\SGARefLink{sga3:iii:corollary:0-7}{0.7} 和
\SGARefLink{sga3:iii:corollary:0-7-bis}{0.7 bis} 的推论,
它们在本篇后文中起重要的技术性作用.}

\medskip
\noindent\SGAInlineRefTarget{sga3:iii:corollary:0-7}\textbf{推论 0.7.}
\textit{设 \(X\) 是一个 \(S\)-概形.}
\begin{enumerate}
 \item[\SGAInlineRefTarget{sga3:iii:corollary:0-7:item:a}(a)] \textit{\(X\) 的每个在
 \(X_{\mathscr J}\) 上诱导恒等态射的 \(S\)-自同态都是自同构.}

 \item[\SGAInlineRefTarget{sga3:iii:corollary:0-7:item:b}(b)] \textit{有群的正合列}
 \[
  0\longrightarrow
  \operatorname{Hom}_{\mathscr O_{X_0}}
    (\Omega^1_{X_0/S_0},\mathscr J\mathscr O_X)
  \xrightarrow{\ i\ }
  \operatorname{Aut}_S(X)
  \longrightarrow
  \operatorname{Aut}_{S_{\mathscr J}}(X_{\mathscr J}).
 \]

 \item[\SGAInlineRefTarget{sga3:iii:corollary:0-7:item:c}(c)] \textit{令
 \(\operatorname{Aut}_S(X)\) 通过结构迁移作用在第一个群上.
 则对每个 \(u\in\operatorname{Aut}_S(X)\) 和
 \(m\in\operatorname{Hom}_{\mathscr O_{X_0}}
 (\Omega^1_{X_0/S_0},\mathscr J\mathscr O_X)\), 有}
 \[
  i(u\cdot m)=u\,i(m)\,u^{-1}.
 \]
\end{enumerate}

% Source page: Exp3 component-local p. 16; combined-reader p. 126; printed p. 116.
\noindent\textit{证明.}
由 \SGARefLink{sga3:iii:proposition:0-2:item:i}{0.2(i)},
\(\operatorname{Hom}_{X^+}(X,X)\) 是
\(\operatorname{Hom}_{X^+}(X,L_X)\) 下的主齐性集,
因为它显然非空: 它包含由 \(X\) 的恒等自同构给出的标定点.
\footnote{编者注: 以下对原文作了展开.}
因此, 映射 \(m\mapsto m\cdot\operatorname{id}_X\) 诱导双射
\[
 \operatorname{Hom}_{\mathscr O_{X_0}}
   (\Omega^1_{X_0/S_0},\mathscr J\mathscr O_X)
 =L_X(X)
 \xrightarrow{\sim}
 \operatorname{Hom}_{X^+}(X,X).
\]

设 \(m\in L_X(X)\), 并设
\(f=m'\cdot\operatorname{id}_X\) 是
\(\operatorname{Hom}_{X^+}(X,X)\) 的一个元素.
将 \SGARefLink{sga3:iii:proposition:0-2:item:ii}{0.2(ii)} 用于 \(f\), 得
\[
 f\circ(m\cdot\operatorname{id}_X)
 =L_f(X)(m)\cdot f
 =L_f(X)(m)\cdot(m'\cdot\operatorname{id}_X).
\]
另一方面, 由于 \(f\) 是 \(X\) 的一个 \(X^+\)-自同态,
有 \(f_{\mathscr J}=\operatorname{id}_{X_{\mathscr J}}\),
从而 \(f_0=\operatorname{id}_{X_0}\).
由于 \(L_f\) 仅依赖于 \(f_0\)
（参见 \SGARefLink{sga3:iii:corollary:0-1-10}{0.1.10} 中的编者注 (17)）,
可得 \(L_f(X)(m)=m\). 因而前面的等式成为
\[
 (m'\cdot\operatorname{id}_X)\circ
 (m\cdot\operatorname{id}_X)
 =m\cdot(m'\cdot\operatorname{id}_X)
 =(m+m')\cdot\operatorname{id}_X.
\]
这表明双射 \(m\mapsto m\cdot\operatorname{id}_X\)
将 \(\operatorname{Hom}_{X^+}(X,L_X)\) 的群运算对应到
\(X\) 的 \(X^+\)-自同态的复合.

首先由此可得, \(\operatorname{Hom}_{X^+}(X,X)\) 的每个元素均可逆,
这证明了~\SGARefLink{sga3:iii:corollary:0-7:item:a}{(a)};
其次可得正合列
\[
 0\longrightarrow \operatorname{Hom}_{X^+}(X,L_X)
 \xrightarrow{\ i\ }\operatorname{Aut}_S(X)
 \longrightarrow\operatorname{Aut}_{S_{\mathscr J}}(X_{\mathscr J}),
\]
这证明了~\SGARefLink{sga3:iii:corollary:0-7:item:b}{(b)}.

刚刚定义的态射 \(i\) 关于 \(X\) 的同构具有函子性,
因为它通过结构由 \(L_X\) 在 \(X^+\) 上对 \(X\) 的作用定义,
而由\SGARefLink{sga3:iii:proposition:0-2:item:ii}{命题~0.2(ii)},
这一作用本身就关于 \(X\) 具有函子性.
\footnote{编者注: 以下对原文作了展开.}
因此, \(X\) 在 \(S\) 上的每个自同构 \(u\) 都通过结构迁移诱导同构
\[
 h:\operatorname{Hom}_{X^+}(X,L_X)
   \xrightarrow{\sim}\operatorname{Hom}_{X^+}(X,L_X)
 \quad\text{及}\quad
 f:\operatorname{Aut}_S(X)
   \xrightarrow{\sim}\operatorname{Aut}_S(X)
\]
使下图交换.

\SGAThreeNativeDiagram[0.70\textwidth]
  {figures/exp3/diagram_022_exp3_p016_cor07_functoriality_square.tex}
  {\(i\) 关于结构迁移的函子性.}
  {fig:exp3-cor07-functoriality}
  {Polo--Gille Exp3-14oct24.pdf, 分篇页 16; 合订本页 126; 印刷页 116}

因此 \(f\circ i=i\circ h\). 另一方面, \(f\) 可由下列交换图描述.

\SGAThreeNativeDiagram[0.40\textwidth]
  {figures/exp3/diagram_023_exp3_p016_cor07_conjugation_square.tex}
  {自同构 \(f\) 通过关于 \(u\) 的共轭作用.}
  {fig:exp3-cor07-conjugation}
  {Polo--Gille Exp3-14oct24.pdf, 分篇页 16; 合订本页 126; 印刷页 116}

% Source page: Exp3 component-local p. 17; combined-reader p. 127; printed p. 117.
换言之,
\[
 f(a)=u\circ a\circ u^{-1}
 \qquad\text{对每个 }a\in\operatorname{Aut}_S(X).
\]
现在写出 \(i(h(m))=f(i(m))\), 就得到所述公式.

% Continue with Corollary 0.7 bis.
